\documentclass[varwidth=14cm,border=10pt]{standalone}
\usepackage{amsmath,xcolor,array}
\usepackage{fontspec,unicode-math}
\setmainfont{DejaVu Sans}[ItalicFont={DejaVu Sans},ItalicFeatures={FakeSlant=0.2},BoldItalicFont={DejaVu Sans Bold},BoldItalicFeatures={FakeSlant=0.2}]
\setmathfont{texgyredejavu-math.otf}
\setmathfont{DejaVu Sans}[range={up,bfup}]
\setmathfont{DejaVu Sans}[range={it,bfit},FakeSlant=0.2]
\setmathfont{texgyredejavu-math.otf}[range={"2202,"1D715,"03B5,"03F5,"2112,"2131}]
\pagecolor[HTML]{1E1E2E}\color[HTML]{F4F4F4}
\definecolor{accent}{HTML}{89B4FA}
\definecolor{warn}{HTML}{F38BA8}
\definecolor{good}{HTML}{A6E3A1}
\newcommand{\heading}[1]{\par\medskip{\large\bfseries\color{accent}#1}\par\smallskip}
\setlength{\parindent}{0pt}
\raggedright
\begin{document}

{\Large\bfseries L04 — Constrained optimisation: KKT, multipliers, SQP}

\heading{1. The problem, and one rule about writing it}
\[
  \min_{x\in\mathbb{R}^n} f(x)
  \quad\text{s.t.}\quad
  g_j(x)\le 0,\ \ j=1..m, \qquad h_k(x)=0,\ \ k=1..p
\]
$g_j$ is {\bfseries active} where $g_j=0$, {\bfseries inactive} where $g_j<0$.
An equality is always active.

{\color{warn}Always write $g\le0$, and non-dimensionalise.} For
$\sigma\le\sigma_{\text{allow}}$ write
\[
  g=\frac{\sigma}{\sigma_{\text{allow}}}-1\le0,
  \qquad\text{not}\qquad \sigma-\sigma_{\text{allow}}\le0 .
\]
Both are correct; only the first makes the multipliers {\itshape comparable}.
The reference must be a {\bfseries constant}.

\heading{2. Why a multiplier has to exist}
Suppose $x^\star$ minimises $f$ on $g=0$ but $\nabla f$ is {\itshape not}
parallel to $\nabla g$. Then $-\nabla f$ has a component {\itshape along} the
constraint surface. Step that way: you stay feasible and $f$ falls --- so
$x^\star$ was not optimal. Therefore at an optimum
\[
  \nabla f(x^\star)\parallel\nabla g(x^\star)
  \qquad\Longrightarrow\qquad
  \nabla f(x^\star)=-\lambda\,\nabla g(x^\star).
\]
{\bfseries The gradient does not vanish at a constrained optimum --- it is
held.} The fence pushes back exactly as hard as the slope pulls.

\heading{3. Equalities: the Lagrangian}
\[
  \mathcal{L}(x,\beta)=f(x)+\sum_{k=1}^{p}\beta_k h_k(x),
  \qquad
  \frac{\partial\mathcal{L}}{\partial x_i}=0,\quad
  \frac{\partial\mathcal{L}}{\partial \beta_k}=0 .
\]
$n+p$ equations, $n+p$ unknowns. The $\beta$ derivative just returns $h_k=0$.

\heading{4. Inequalities: the KKT conditions}
With $\mathcal{L}=f+\sum_j\lambda_jg_j$, at a constrained minimum:
\[
  \begin{array}{@{}l@{\qquad}l@{}}
    \nabla f(x^\star)+\sum_j \lambda_j\nabla g_j(x^\star)=0 & \text{stationarity}\\[2pt]
    g_j(x^\star)\le0 & \text{feasibility}\\[2pt]
    \lambda_j\ge0 & \text{sign}\\[2pt]
    \lambda_j\,g_j(x^\star)=0 & \text{complementary slackness}
  \end{array}
\]
Complementary slackness is the useful one: {\bfseries either the constraint is
tight, or its price is zero.} Carry a second set $\sum_k\beta_k\nabla h_k$ for
equalities --- which is why \texttt{fmincon} returns
\texttt{lambda.ineqnonlin} and \texttt{lambda.eqnonlin} separately.

\heading{5. What KKT is for}
KKT is {\bfseries necessary}. If the problem is {\bfseries convex}, KKT is also
{\bfseries sufficient} --- and then the minimum is {\bfseries global}.
Necessary $\Rightarrow$ how we {\itshape search}. Sufficient $\Rightarrow$ how
we {\itshape confirm}.

\heading{6. The multiplier is a price}
\[
  \boxed{\ \frac{df^\star}{db}=-\lambda\ }
\]
Sweeping the bound in $g=x_1+2x_2-b\le0$:
\[
  \begin{array}{@{}r@{\qquad}r@{\qquad}r@{\qquad}l@{}}
    b & f^\star & \lambda & \\
    3.00 & 1.80000 & 1.20000 & \text{active --- and expensive}\\
    4.75 & 0.31250 & 0.50000 & \text{active}\\
    5.50 & 0.05000 & 0.20000 & \text{active --- nearly worthless}\\
    6.00 & 0.00000 & 0.00000 & \text{inactive: the free optimum became feasible}
  \end{array}
\]
{\color{good}A constraint you are not up against is worth nothing.}

\heading{7. KKT turns constrained into unconstrained}
Differentiate $\mathcal{L}=f+\lambda g$ in {\itshape all} arguments, treating
$\lambda$ as a variable like any other:
\[
  \nabla\mathcal{L}=
  \begin{bmatrix}\partial\mathcal{L}/\partial x_1\\ \partial\mathcal{L}/\partial x_2\\ \partial\mathcal{L}/\partial\lambda\end{bmatrix}
  =\begin{bmatrix}2x_1-4+\lambda\\ 2x_2-4+2\lambda\\ x_1+2x_2-4\end{bmatrix}
\]
The first rows are stationarity; the last is $\partial\mathcal{L}/\partial\lambda=g$,
i.e.\ the active constraint. So {\bfseries $\nabla\mathcal{L}=0$ {\itshape is} the
KKT system}: $n+m$ equations in $n+m$ unknowns.

\heading{8. Therefore Newton is unchanged}
Put $\mathcal{L}$ where $f$ was and $(x,\lambda)$ where $x$ was:
\[
  \begin{bmatrix}x\\ \lambda\end{bmatrix}^{(k+1)}
  =\begin{bmatrix}x\\ \lambda\end{bmatrix}^{(k)}-H_{\mathcal{L}}^{-1}\nabla\mathcal{L}
\]
Same formula as Lecture 3. {\itshape Linear} KKT system $\Rightarrow$
$H_{\mathcal{L}}$ constant $\Rightarrow$ converges in {\bfseries one step}.
{\color{warn}Do not generalise that.} The tin can's KKT system is nonlinear, so
$H_{\mathcal{L}}$ is rebuilt each step --- $\|\nabla\mathcal{L}\|$ then runs
$5.1\!\times\!10^{1}\to4.8\to1.1\to8.0\!\times\!10^{-2}\to1.4\!\times\!10^{-4}
\to2.4\!\times\!10^{-10}$: roughly squaring, i.e.\ {\bfseries quadratic
convergence}.

\heading{9. What makes it SQP}
Raw Newton always takes $\alpha=1$ and can diverge from a poor start. SQP keeps
the same direction $d_k=-H_{\mathcal{L}}^{-1}\nabla\mathcal{L}$ but line-searches
$\alpha_k$ on a {\bfseries merit function}
\[
  \Psi(x)=f(x)+\sum_j\lambda_j\max[0,g_j(x)]+\sum_k\lambda_{m+k}|h_k(x)| ,
\]
then minimises $\varphi(\alpha)=\Psi(x^{(k)}+\alpha d_k)$ --- Lecture 3's line
search with $\Psi$ in place of $f$. {\bfseries KKT is a test; SQP is a method.}

\heading{10. The engineering payoff (heat exchanger)}
Every $g$ divided by a fixed reference of its own kind, so all six $\lambda$ are
in one currency and can be {\bfseries ranked}:
\[
  \begin{array}{@{}lr@{}}
    \text{sizing 3} & 10\,000.0 \ \leftarrow\\
    \text{balance 2} & 5210.7\\
    \text{balance 3} & 5110.0\\
    \text{sizing 2} & 4789.3\\
    \text{sizing 1} & 2825.3\\
    \text{balance 1} & 1964.0
  \end{array}
\]
All six active; stage 3's sizing leads at about twice anything else, so
relaxing it 10\% saves roughly $1000\ \mathrm{m}^2$. {\bfseries The ranking is
the answer} --- it tells the engineer which limit to argue about.

\heading{11. Take away}
\begin{enumerate}\itemsep0pt
\item At a constrained optimum the gradient is {\itshape held}, not zero:
      $\nabla f=-\sum\lambda_j\nabla g_j$.
\item The multiplier is a price, $df^\star/db=-\lambda$; a slack constraint is
      worth nothing.
\item KKT is a test, SQP is a method --- and SQP is Lecture 3's Newton step on
      $\mathcal{L}$ in $(x,\lambda)$. No new machinery.
\item Multipliers in different units cannot be compared. Non-dimensionalise
      everything.
\item Report the {\bfseries design}, not only the objective --- and the useful
      information came from the multipliers.
\end{enumerate}

\end{document}
